\section{VGA text buffer 的特殊处理}

\begin{thinkbox}
VGA text buffer 位于物理地址 \hex{B8000}。拆除 identity mapping 后，
通过 \hex{B8000} 访问 VGA 会触发 Page Fault。
内核应该如何继续输出到屏幕？

\textbf{方案分析}：
\begin{itemize}
  \item 方案 A：通过 high-half 直接映射访问 \hex{C00B8000}
        （$= \hex{B8000} + \hex{C0000000}$）
  \item 方案 B：放弃 VGA，只使用串口（serial port）输出
\end{itemize}

我们的内核选择了方案 B：\texttt{printk} 完全通过串口（I/O 端口 \hex{3F8}）输出。
串口使用 \texttt{in}/\texttt{out} 指令访问 I/O 端口空间，不经过内存映射，
因此不受 identity mapping 拆除的影响。

如果将来需要 VGA 输出，只需将 VGA 基地址改为 \hex{C00B8000} 即可。
这正是直接映射区的方便之处——任何物理设备地址加上 \hex{C0000000} 就能访问。
\end{thinkbox}

% ============================================================================

